Fix a transform \(M\) represented in the class.  The preceding lemma supplies a zero in \(|z|\le R_0\), hence every circle of radius larger than \(R_0\) encloses at least one zero.  We first verify the uniform zero count used by the construction.  From the Luxemburg transform bound used in \(\text{lem:zero-localization}\),
\[
 \log\max_{|z|=R_{\mathrm{out}}}|M(z)|
 \le4\psi_{\eta}^2R_{\mathrm{out}}^2.
\]
Jensen's formula, with \(M(0)=1\), therefore shows that the number of zeros in \(|z|\le R_{\mathrm{fac}}\), counted with multiplicity, is at most
\[
 \frac{4\psi_{\eta}^2R_{\mathrm{out}}^2}
 {\log(R_{\mathrm{out}}/R_{\mathrm{fac}})}
 =Q_{\star}\le N_{\mathrm{cert}}. \tag{1}
\]

The grid \(\mathcal Q_{m_{\star}}\) is finite, lies in \((R_0,R_1)\), and has more than \(N_{\mathrm{cert}}\) points separated by \(d_{\star}\).  A zero can have modulus within \(h_{\star}=d_{\star}/3\) of at most one grid point.  Hence (1) and the pigeonhole principle give a grid radius \(\rho_j\) whose circle is at radial distance at least \(h_{\star}\) from every zero in \(|z|le R_{\mathrm{fac}}\).  This circle encloses the zero supplied above, so \(N_{C_j}(M)\ge1\).

Factor the at most \(N_{\mathrm{cert}}\) zeros in the disk by the corresponding disk Blaschke factors and call the remaining zero-free factor \(U\), normalized at the origin.  On \(|z|\le R_1\), the separation just obtained bounds the product of the Blaschke moduli below by
\[
 \left(\frac{h_{\star}}{R_{\mathrm{fac}}+R_1}\right)^{N_{\mathrm{cert}}}.
\]
Applying the Poisson formula to the positive harmonic majorant of \(-\log|U|\), together with the transform upper bound on \(|z|=R_{\mathrm{fac}}\), gives the Harnack factor
\[
 \inf_{|z|\le R_1}|U(z)|
 \ge \exp\{-4(\Lambda_{\star}-1)psi_{\eta}^2R_{\mathrm{fac}}^2\}.
\]
Multiplication gives \(\inf_{C_j}|M|\ge a_{\mathrm{an}}\ge a_{\star}>0\).  All searches are over finite integer or rational sets, and outward interval refinement terminates because the target inequalities are strict.  Thus \(J\), the radii, and the positive dyadic certificate are computable from the certified-real names of the displayed primitive constants, do not use the unknown law, and are common to every \(n\).  This proves every assertion.